% !PRINTED 628
\weberpage{628}
\begin{center}
\S.\,173.

\gesperrt{Darstellung der Primfactoren von $p$.}
\end{center}

Zur Darstellung der Primfunctionale des Körpers $\Omega_m$ können wir ein Verfahren anwenden, was wir im \S.\,157 kennen gelernt haben, welches sich hier in Folge des Umstandes, dass
\[ 1, r, r^2, \ldots, r^{\mu-1} \]
eine Basis von $\mathfrak{o}$ ist, wesentlich vereinfacht.

Die natürliche Primzahl $q$ ist, wie wir schon gesehen haben, die $m^{\mathrm{te}}$ Potenz einer im Körper $\Omega_m$ existirenden Primzahl $\sigma$; mit dieser brauchen wir uns nicht weiter zu beschäftigen, und betrachten daher hier nur die von $q$ verschiedenen Primzahlen $p$.
% [sic] Source prints the m-th power of sigma here; the ramification index of the prime above q is mu = phi(m). Preserve the printed m.
Es sei $\mathfrak{p}$ ein Primfactor von $p$, und $\mathfrak{A}$ habe dieselbe Bedeutung wie oben. Wenn $p$ zum Exponenten $f$ gehört und
\[ ef = \varphi(m) = \mu \]
ist, so besteht die Gruppe $\mathfrak{A}$ aus den Zahlen
\[ 1, p, p^2, \ldots, p^{f-1}. \]
Wenn nun
\[ f(t) = N(t-r) = \prod^n (t - r^n) \]
das Polynom von $t$ vom Grade $\mu$ bedeutet, dessen Wurzeln die $\mu$ Grössen $r^n$ sind, so können wir dies so in Factoren zerlegen, dass jeder Factor einer der Nebengruppen von $\mathfrak{A}$ entspricht. Setzen wir nämlich
\weq{1}{F_1(t) = \prod_{0,\,f-1}^{h} (t - r^{p^h}),}
und allgemein für jedes durch $q$ nicht theilbare $n$
\weq{2}{F_n(t) = \prod_{0,\,f-1}^{h} (t - r^{np^h}),}
so ist $F_n(t)$ mit $F_{n'}(t)$ identisch, wenn $n$ und $n'$ in dieselbe Nebengruppe $\mathfrak{A}\xi_i$ gehören. Wenn wir also mit $\xi_1, \xi_2, \ldots, \xi_e$ das im \S.\,172, (5) angewandte Zahlensystem verstehen, so ist
\weq{3}{f(t) = F_{\xi_1}(t)\, F_{\xi_2}(t) \ldots F_{\xi_e}(t).}

Nun ist aber nach dem binomischen Lehrsatze:
\[ (t - r^{p^h})^p \equiv (t^p - r^{p^{h+1}}) \, (\mathrm{mod}\ p), \]
%CONT>>
